% Source printed pp. 195-204 / PDF pp. 203-212.
\section{代数空间与概形}\label{src:5.5}
\begin{epigraph}
但谁能量化\\
空间的代数,\\
或称量那些浮游的世界,\\
各在其位?\\
谁能胜过黑暗?\\
又有哪台计算机知道\\
美如何诞生---\\
贝壳, 星辰与玫瑰?
\epigraphauthor{Jean Kenward, \textit{技术人员}}
\end{epigraph}
我们证明若干结果, 给出代数空间成为概形的条件:
\begin{itemize}
\item 拟分离代数空间在一个稠密开子空间上是概形 (定理~\bookref{5.5.1});
\item 代数空间的 Zariski 主定理: 可表出, 拟有限且分离的态射可分解为开浸入与有限态射的复合 (定理~\bookref{5.5.9}). 它有一些很好的推论:
\begin{itemize}
\item 在概形上局部拟有限且分离的代数空间是概形 (推论~\bookref{5.5.7}),
\item 若 Deligne--Mumford 叠的对角线分离且拟紧, 则对角线拟仿射 (推论~\bookref{5.5.8}), 特别地可由概形表出, 且
\item 稳定子平凡的代数叠是代数空间 (定理~\bookref{5.5.10}), 推广了定理~\bookref{4.6.6};
\end{itemize}
\item 代数空间的 Serre 与 Chevalley 仿射性判据 (定理~\bookref{5.5.16} 和~\bookref{5.5.20});
\item 若 $X$ 是域 $k$ 上局部有限型的拟分离代数空间, 且 $X_{\overline k}$ 的每个有限点集都包含在某个仿射开集中 (例如 $X_{\overline k}$ 拟射影), 则 $X$ 是概形 (命题~\bookref{5.5.27});
\item 域上局部有限型的拟分离群代数空间是概形 (定理~\bookref{5.5.28}); 且
\item 分离的一维代数空间是概形 (定理~\bookref{5.5.32}).
\end{itemize}

\subsection{代数空间在稠密开集上是概形}\label{src:5.5.1-subsection}
\begin{theorem}\label{src:5.5.1}
每个拟分离代数空间都有一个是概形的稠密开子空间.
\end{theorem}
\begin{proof}
可假设 $X$ 是拟紧拟分离的代数空间. 设 $f:V\to X$ 是平展呈示, 其中 $V$ 是仿射概形. 由于 $X$ 拟分离, $f:V\to X$ 拟紧, 且存在开代数子空间 $U\subseteq X$, 使得 $f^{-1}(U)\to U$ 有限. 由习题~\bookref{5.3.2}, $U$ 同构于有限抽象群 $G$ 在概形 $V$ 上自由作用所得的商叠 $[V/G]$. 若 $V_1\subseteq V$ 是稠密仿射开子概形, 则 $V_2=\bigcap_{g\in G}gV_1$ 是 $V$ 的 $G$-不变拟仿射开子概形, 特别地是分离的. 重复此论证, 可选取稠密仿射开子概形 $V_3\subseteq V_2$, 此时 $V_4=\bigcap_{g\in G}gV_3$ 是 $G$-不变仿射开子概形. 命题~\bookref{5.3.3} 蕴含 $V_4/G\cong\Spec A^G$ 是 $U$ 的稠密仿射开代数子空间. 另见~\cite[II.6.7]{Knu71} 和~\cite[Tag 06NN]{SP}.
\end{proof}
\begin{remark}\label{src:5.5.2}
若 $X$ 不拟分离, 上述结果未必成立, 例如 $\mathbb A^1/\mathbb Z$ (例~\bookref{4.9.36}).
\end{remark}
\begin{corollary}\label{src:5.5.3}
整的拟分离代数空间有良定义的分式域.
\end{corollary}
\begin{exercise}\label{src:5.5.4}
设 $G$ 是作用于拟紧拟分离代数空间 $U$ 的有限抽象群. 证明 $U$ 有稠密的 $G$-不变仿射开子概形.
\end{exercise}

\subsection{Zariski 主定理}\label{src:5.5.2-subsection}
我们证明 Deligne--Mumford 叠的 Zariski 主定理 (\bookref{5.5.9}), 它的一个重要应用是, 代数空间的局部拟有限且分离的态射可由概形表出 (推论~\bookref{5.5.7}). 证明依赖拟凝聚层理论, 特别是拟紧拟分离态射的分解
\[
f:\mathcal X\longrightarrow\Spec_{\mathcal Y}f_*\mathcal O_{\mathcal X}\longrightarrow\mathcal Y.
\]
代数空间的拟有限态射的定义见定义~\bookref{4.3.34}, 概形的 Zariski 主定理的讨论见~\S\bookref{A.7}.
\begin{proposition}\label{src:5.5.5}
Deligne--Mumford 叠的可表出, 拟有限且分离的态射 $f:\mathcal X\to\mathcal Y$ 可分解为开浸入 $\mathcal X\hookrightarrow\Spec_{\mathcal Y}f_*\mathcal O_{\mathcal X}$ 与仿射态射 $\Spec_{\mathcal Y}f_*\mathcal O_{\mathcal X}\to\mathcal Y$ 的复合. 特别地 $f$ 拟仿射.
\end{proposition}
\begin{proof}
由于 $f_*\mathcal O_{\mathcal X}$ 的构造与 $\mathcal Y$ 上的平坦基变换相容 (习题~\bookref{5.1.43}), 分解 $f:\mathcal X\to\Spec_{\mathcal Y}f_*\mathcal O_{\mathcal X}\to\mathcal Y$ 的形成也如此. 因而此陈述在 $\mathcal Y$ 上是平展局部的. 特别地, 可假设 $\mathcal Y=Y$ 是仿射概形, $\mathcal X=X$ 是代数空间. 将 $Y$ 换成 $\Spec_Y f_*\mathcal O_X$ 后, 可假设 $f_*\mathcal O_X=\mathcal O_Y$, 须证明 $f:X\to Y$ 是开浸入.

由于 $X$ 拟紧, 存在来自仿射概形的平展呈示 $\pi:U\to X$. 由于 $X$ 分离, $U\to X$ 也分离. 复合
\[
U\xrightarrow{\pi}X\xrightarrow{f}Y
\]
是概形的拟有限态射, 因而可在每个点 $y\in Y$ 附近应用拟有限态射的平展局部化 (\bookref{A.7.1}): 将 $Y$ 换成平展邻域后, 可假设 $U=U_1\amalg U_2$, 其中 $U_1\to Y$ 有限, $(U_2)_y=\varnothing$. 则 $\pi(U_1)$ 既开 (因为 $\pi$ 平展) 又闭 (因为 $U_1\to Y$ 有限, $X\to Y$ 分离). 因此 $X=X_1\amalg X_2$, 其中 $X_1=\pi(U_1)$, $(X_2)_y=\varnothing$. 这说明 $\mathcal O_Y=f_*\mathcal O_X$ 是拟凝聚 $\mathcal O_X$-代数的乘积 $\mathcal A_1\times\mathcal A_2$, 从而也可分解 $Y=Y_1\amalg Y_2$, 使得 $y\in Y_1$, 且对 $i=1,2$ 有 $f(Y_i)\subseteq X_i$. 将 $Y$ 换成 $Y_1$ 后, 复合 $U\to X\to Y$ 有限, 引理~\bookref{5.5.6} 蕴含 $X$ 仿射. 因此 $X=Y=\Spec_Y f_*\mathcal O_X$.
\end{proof}
\begin{lemma}\label{src:5.5.6}
设 $U\to X$ 是代数空间的满射平展态射, $X\to Y$ 是代数空间的分离态射. 若复合 $U\to X\to Y$ 有限, 则 $X\to Y$ 也有限.
\end{lemma}
\begin{proof}
此陈述在 $Y$ 上是平展局部的, 故可假设 $Y$ 与 $U$ 仿射. 由于 $X\to Y$ 分离, $U\to X$ 也有限. $X$ 可与仿射概形的有限平展等价关系 $R:=U\times_X U\rightrightarrows U$ 的商 $U/R$ 等同, 因而命题~\bookref{5.3.3} 蕴含 $X$ 仿射. 由于 $U\to Y$ 固有, $X\to Y$ 也固有. $X\to Y$ 是概形的固有拟有限态射, 因而有限 (推论~\bookref{A.7.5}).
\end{proof}
\begin{corollary}\label{src:5.5.7}
代数空间的局部拟有限且分离的态射可由概形表出. 特别地, 在概形上局部拟有限且分离的代数空间是概形.
\end{corollary}
\begin{proof}
只需说明, 若 $X\to Y=\Spec A$ 局部拟有限且分离, 则 $X$ 是概形. 由于为概形是 Zariski 局部的性质, 可假设 $X$ 拟紧. 此时可应用命题~\bookref{5.5.5}.
\end{proof}
\begin{corollary}\label{src:5.5.8}
对角线分离且拟紧的 Deligne--Mumford 叠的对角线拟仿射. 特别地, 拟分离代数空间有拟仿射对角线.
\end{corollary}
\begin{proof}
由于对角线可表出, 拟有限且分离, 可应用命题~\bookref{5.5.5}. 注意代数空间的对角线是单态射, 因而分离.
\end{proof}
与概形的情形一样, 可细化命题~\bookref{5.5.5}, 得到 Zariski 主定理的一个版本.
\begin{theorem}[Zariski 主定理]\label{src:5.5.9}
Noether Deligne--Mumford 叠的可表出, 拟有限且分离的态射 $f:\mathcal X\to\mathcal Y$ 可分解为稠密开浸入 $\mathcal X\hookrightarrow\widetilde{\mathcal Y}$ 与有限态射 $\widetilde{\mathcal Y}\to\mathcal X$ 的复合.
\end{theorem}
\begin{proof}
设 $\mathcal A\subseteq f_*\mathcal O_{\mathcal X}$ 是 $\mathcal O_{\mathcal Y}\to f_*\mathcal O_{\mathcal X}$ 的整闭包: 对来自概形的平展态射 $T\to\mathcal Y$, $\mathcal A(T\to\mathcal Y)$ 是 $\Gamma(T,\mathcal O_T)\to\Gamma(\mathcal X_T,\mathcal O_{\mathcal X_T})$ 的整闭包. 由于整闭包与平展扩张相容 (命题~\bookref{A.7.4}), $\mathcal A$ 是拟凝聚 $\mathcal O_{\mathcal Y}$-代数层. 由习题~\bookref{5.1.33}, 可写 $\mathcal A=\operatorname{colim}\mathcal A_\lambda$, 为有限型 $\mathcal O_{\mathcal Y}$-代数的余极限. 由于 $\mathcal Y$ 拟紧, 存在来自仿射概形的平展呈示 $p:U\to\mathcal Y$. 基变换 $\mathcal X_U\to U$ 是代数空间的拟有限分离态射, 因而由推论~\bookref{5.5.7} 它是概形的态射. $p^*\mathcal A=\operatorname{colim}p^*\mathcal A_\lambda$ 可与 $\mathcal O_U\to f_{U,*}\mathcal O_{\mathcal X_U}$ 的整闭包等同, 因而由概形的 Zariski 主定理 (\bookref{A.7}), 对 $\lambda\gg0$, $\mathcal X_U\to\Spec_U p^*\mathcal A_\lambda$ 是开浸入, $\Spec_U p^*\mathcal A_\lambda\to U$ 有限. 由平展下降, 对 $\lambda\gg0$, $\mathcal X\to\Spec_{\mathcal Y}\mathcal A_\lambda$ 也是开浸入, $\Spec_{\mathcal Y}\mathcal A_\lambda\to\mathcal Y$ 有限. 另见~\cite[II.6.15]{Knu71}, \cite[定理 A.2]{LMB00}, \cite[定理 $7.2.10$]{Ols16} 和~\cite[Tag 05W7]{SP}.
\end{proof}

\subsection{代数空间的刻画}\label{src:5.5.3-subsection}
现在可去掉定理~\bookref{4.6.6} 和推论~\bookref{4.6.9} 中关于对角线可由概形表出的假设.
\begin{theorem}[代数空间的刻画 II]\label{src:5.5.10}
对代数叠 $\mathcal X$, 以下条件等价:
\begin{enumerate}[(1)]
\item 叠 $\mathcal X$ 是代数空间,
\item 对角线 $\mathcal X\to\mathcal X\times\mathcal X$ 是单态射, 且
\item $\mathcal X$ 的每个点都有平凡稳定子.
\end{enumerate}
\end{theorem}
\begin{proof}
只需说明 (2)$\Rightarrow$(1). 由于对角线 $\Delta:\mathcal X\to\mathcal X\times\mathcal X$ 是单态射, 它局部拟有限且分离. 将证明 $\Delta$ 可由概形表出, 从而可应用定理~\bookref{4.6.6}. 为此, 设 $(a,b):T\to\mathcal X\times\mathcal X$ 是来自概形的态射, 并考虑沿对角线的基变换 $\Isom_T(a,b)\to T$. 由于 $\Isom_T(a,b)$ 在 $T$ 上局部拟有限且分离, 由推论~\bookref{5.5.7} 它是概形.
\end{proof}
\begin{corollary}[可表态射的刻画 II]\label{src:5.5.11}
代数叠的态射 $\mathcal X\to\mathcal Y$ 可表出, 当且仅当对每个几何点 $x\in\mathcal X(k)$, 自同构群上的映射 $G_x\to G_{f(x)}$ 单射.
\end{corollary}
若 $X$ 是层, 且存在来自概形的可由概形表出的满射平展态射 $U\to X$, 则 $X$ 是代数空间 (由定义). 事实上, 若只要求 $U$ 是代数空间, $U\to X$ 光滑且可表出 (由代数空间表出), $X$ 仍是代数空间. 此外, 概形的平展等价关系的商的存在性 (定理~\bookref{4.4.13}(2)) 可推广到代数空间的平展等价关系.
\begin{corollary}\label{src:5.5.12}
\begin{enumerate}[(1)]
\item 若 $X$ 是 $\Sch_{\text{ét}}$ 上的层, 且存在来自代数空间的满射, 平展 (分别为光滑) 且可表出的态射 $U\to X$, 则 $X$ 是代数空间.
\item 若 $R\rightrightarrows U$ 是代数空间的平展 (分别为光滑) 等价关系, 则商 $U/R$ 是代数空间.
\end{enumerate}
\end{corollary}
\begin{remark}\label{src:5.5.13}
将``平展''换成``fppf''后, 上述陈述仍成立; 见定理~\bookref{7.3.1} 和推论~\bookref{7.3.6}.
\end{remark}
\begin{proof}[推论~\bookref{5.5.12} 的证明]
由定理~\bookref{4.4.13}(1), (1) 中的层 $X$ 与 (2) 中的商层 $U/R$ 都是代数叠, 结论由定理~\bookref{5.5.10} 推出.
\end{proof}
\begin{corollary}\label{src:5.5.14}
代数空间的固有拟有限态射 (分别为固有单态射) 有限 (分别为闭浸入).
\end{corollary}
\begin{proof}
固有拟有限态射可由概形表出. 因而此陈述由概形的对应结果 (推论~\bookref{A.7.5}) 和平展下降推出.
\end{proof}
\begin{exercise}\label{src:5.5.15}
证明 $\Sch_{\text{ét}}$ 上的预叠 $\underline{\mathrm{AlgSp}}$ 是叠, 其中在概形 $T$ 上的对象为 $T$ 上的代数空间, 态射对应于代数空间的 Cartesian 图.
\end{exercise}

\subsection{仿射性判据}\label{src:5.5.4-subsection}
\begin{theorem}[Serre 仿射性判据]\label{src:5.5.16}
设 $X$ 是拟紧拟分离 (分别为 Noether) 代数空间. 若函子 $\Gamma(X,-)$ 在拟凝聚层 (分别为凝聚层) 范畴上正合, 则 $X$ 是仿射概形.
\end{theorem}
\begin{proof}
将证明典范态射 $\pi:X\to Y:=\Spec\Gamma(X,\mathcal O_X)$ 是固有单态射. 由于固有单态射是闭浸入 (推论~\bookref{5.5.14}), 这将蕴含 $X$ 仿射. 证明受到良模空间的性质的启发 (见~\S\bookref{7.5}).

\textit{断言:} 若 $g:Y'\to Y$ 是代数空间的态射, 则基变换 $\pi':X':=X\times_Y Y'\to Y'$ 有下列性质:
\begin{enumerate}[(a)]
\item $\pi'_*$ 诱导 $X'$ 与 $Y'$ 上的拟凝聚层范畴的等价.
\item $\mathcal O_{Y'}\to\pi'_*\mathcal O_{X'}$ 是同构.
\item $X'\to Y'$ 是同胚.
\end{enumerate}
由平坦基变换 (习题~\bookref{5.1.43}), 性质 (a) 与 (b) 在 $Y'$ 上是平展局部的, 故可假设 $Y'=\Spec B$. 将说明对 $G\in\QCoh(Y')$ 和 $F\in\QCoh(X')$, 伴随态射 $G\to\pi'_*\pi'^*G$ 与 $\pi'^*\pi'_*F\to F$ 是同构. 这立即给出 (a), 也给出 (b) (取 $G=\mathcal O_{Y'}$). 对第一个伴随态射, 为作为 $\mathcal O_Y$-模的 $G$ 选取自由呈示 $\mathcal O_Y^{\oplus J}\to\mathcal O_Y^{\oplus I}\to g_*G\to0$. 由于 $\pi_*$ 正合, 有右正合列的态射
\[
\begin{tikzcd}[column sep=small]
\mathcal O_Y^{\oplus J}\arrow[r]\arrow[d]&\mathcal O_Y^{\oplus I}\arrow[r]\arrow[d]&g_*G\arrow[r]\arrow[d]&0\\
\pi_*\pi^*(\mathcal O_Y^{\oplus J})\arrow[r]&\pi_*\pi^*(\mathcal O_Y^{\oplus I})\arrow[r]&\pi_*\pi^*g_*G\arrow[r]&0.
\end{tikzcd}
\]
由于 $\pi_*\mathcal O_X=\mathcal O_Y$, 左侧两条竖直箭头是同构. 因此 $g_*G\to\pi_*\pi^*g_*G\cong g_*\pi'_*\pi'^*G$ 是同构. 由于 $g_*$ 忠实正合, $G\to\pi'_*\pi'^*G$ 也是同构. 为说明 $\pi'^*\pi'_*F\to F$ 是同构, 设 $K$ 与 $Q$ 为 $\pi'^*\pi'_*F\to F$ 的核与余核. 由于 $\pi_*$ 正合, $g_*$ 忠实正合, $\pi'_*$ 正合. 由于 $\pi'_*\pi'^*\pi'_*F\to\pi'_*F$ 是同构 (将第一个伴随态射的结论应用于 $\pi'_*F$), 得到 $\pi'_*K=\pi'_*Q=0$. 因而只需说明, 对 $X'$ 上的拟凝聚层 $F'$, $F'\neq0$ 蕴含 $\pi'_*F'\neq0$. 若 $x:\Spec k\to X'$ 是满足 $F\otimes k:=x^*F\neq0$ 的几何点, 沿复合 $\pi'\circ x:\Spec k\to Y'$ 作基变换后, 可假设 $Y'=\Spec k$, 且 $x:\Spec k\to X'$ 是 $\pi'$ 的截面. 由于定义在 $k$ 上的代数空间的每个 $k$-点都是闭点, $x:\Spec k\to X'$ 是闭浸入, 从而 $F\to x_*x^*F=F\otimes k$ 满射. 由 $\pi'_*$ 的正合性, $\pi'_*F\to\pi'_*x_*x^*F=F\otimes k$ 满射, 因而 $\pi'_*F\neq0$.

为说明 (c), 若 $y:\Spec k\to Y$ 是几何点, 则由 (b) 有 $\Gamma(X_y,\mathcal O_{X_y})=k$, 从而纤维 $X_y$ 非空. 另一方面, 若 $x,x'\in X_y(k)$ 是不同点, 它们必为闭点, 则 $\mathcal O_{X_y}\to\mathcal O_{\{x,x'\}}$ 满射. 由于 $\pi_*$ 正合, 还得到满射 $k=\Gamma(X_y,\mathcal O_{X_y})\to k\oplus k$, 矛盾. 为说明 $\pi'$ 闭, 设 $Z\subseteq X'$ 是闭子空间, $q:Z\to\operatorname{im}(Z)$ 是到其概形论像的态射. 则 $\mathcal O_Z\to q_*\mathcal O_{\operatorname{im}(Z)}$ 是同构, 且 $q_*$ 正合. 将上述满射性结果应用于 $q$, 得到 $q$ 满射, 因而 $\pi'(Z)$ 是闭的.

此断言蕴含 $X\to Y$ 普遍闭. 为说明 $X\to Y$ 是单态射, 即对角线 $\Delta:X\to X\times_Y X$ 是同构, 注意 (b) 蕴含 $\mathcal O_{X\times_Y X}\to\Delta_*X$ 沿第一个投影 $p_1:X\times_Y X\to X$ 的推前是同构. 因此将 (a) 应用于 $p_1$, 得到 $\mathcal O_{X\times_Y X}\to\Delta_*X$ 是同构. Zariski 主定理 (\bookref{A.7.3}) 蕴含 $\Delta$ 是开浸入. 将 (c) 应用于 $p_1$, 得到 $p_1:|X\times_Y X|\to|X|$ 是双射. 因而 $\Delta$ 是同构. 还须说明 $X\to Y$ 有限型. 设 $p:U=\Spec A\to X$ 是平展呈示. 由于 $X$ 分离, $R:=U\times_X U$ 是 $U\times_Y U=\Spec(A\otimes_{\Gamma(X,\mathcal O_X)}A)$ 的闭子概形. 因而 $R=\Spec B$ 仿射. 以 $s$ 与 $t$ 记两个映射 $A\rightrightarrows B$, 有交换图
\[
\begin{tikzcd}[row sep=1.5em,column sep=1.7em]
&A\arrow[dr]\arrow[drr,"s"]&&\\
\Gamma(X,\mathcal O_X)\arrow[ur]\arrow[dr]&&A\otimes_{\Gamma(X,\mathcal O_X)}A\arrow[r,two heads]&B.\\
&A\arrow[ur]\arrow[urr,"t"']&&
\end{tikzcd}
\]
由于 $p:U\to X$ 平展, $t:A\to B$ 有限型, 可选取 $A$-代数生成元 $b_1,\ldots,b_n\in B$. 对每个 $i$, 为 $b_i$ 选取原像 $\sum_j a_{ij}\otimes a'_{ij}\in A\otimes_{\Gamma(X,\mathcal O_X)}A$. 注意 $s(a_{ij})$ 也通过 $t$ 将 $B$ 生成为 $A$-代数. 元素 $a_{ij}\in\Gamma(X,p_*\mathcal O_U)=A$ 定义 $\mathcal O_X$-代数的同态 $\mathcal O_X[z_{ij}]\to p_*\mathcal O_U$, 将 $z_{ij}$ 送到 $a_{ij}$. 沿 $p$ 的拉回可与 $\mathcal O_U[z_{ij}]\to p^*p_*\mathcal O_U\cong t_*\mathcal O_R$ 等同, 其中最后的等价来自平坦基变换 (习题~\bookref{5.1.43}); 此映射满射恰是因为 $s(a_{ij})$ 在 $t$ 下生成 $B$. 由平展下降, $\mathcal O_X[z_{ij}]\to p_*\mathcal O_U$ 满射, 因而 $\Gamma(X,\mathcal O_X)[z_{ij}]\to A$ 也满射. 因此 $\Gamma(X,\mathcal O_X)\to A$ 有限型, 由平展下降, $X\to Y$ 也有限型.

最后, 在 Noether 情形, 每个拟凝聚层都是凝聚层的余极限 (习题~\bookref{5.1.28}), 且 $\Gamma(X,-)$ 与余极限相容. 假设 $\Gamma(X,-)$ 在凝聚层上正合. 给定 $X$ 上拟凝聚层的满射 $p:F\twoheadrightarrow G$, 将 $G=\operatorname{colim}_i G_i$ 写成凝聚层的余极限, 并选取满射到 $G_i$ 的凝聚子层 $F_i\subseteq p^{-1}(G_i)$. 则 $\Gamma(X,F_i)\twoheadrightarrow\Gamma(X,G_i)$, 且复合
\[
\operatorname{colim}_i\Gamma(X,F_i)\longrightarrow\Gamma(X,F)\longrightarrow\Gamma(X,G)=\operatorname{colim}_i\Gamma(X,G_i)
\]
满射. 因而 $\Gamma(X,F)\to\Gamma(X,G)$ 满射, 得到 $\Gamma(X,-)$ 在拟凝聚层上正合. 另见~\cite[定理 III.2.5]{Knu71}, \cite[定理 $8.7$]{Ryd15} 和~\cite[Tag 07V6]{SP}.
\end{proof}
\begin{corollary}\label{src:5.5.17}
设 $X$ 是拟紧拟分离 (分别为 Noether) 代数空间. 则 $X$ 是仿射概形, 当且仅当对每个拟凝聚层 (分别为凝聚层) $F$ 和 $i>0$, 都有 $\mathrm H^i(X,F)=0$.
\end{corollary}
\begin{proof}
若 $X$ 仿射, 则定理~\bookref{5.1.40} 给出拟凝聚层上同调的消失. 反之, 拟凝聚层 (分别为凝聚层) 上同调的消失蕴含 $\Gamma(X,-)$ 在拟凝聚层 (分别为凝聚层) 范畴上正合: 若 $0\to F_1\to F_2\to F_3\to0$ 正合, 则由于 $\mathrm H^1(X,F_1)=0$, $\Gamma(X,F_2)\to\Gamma(X,F_3)$ 满射.
\end{proof}
\begin{remark}\label{src:5.5.18}
给定 Deligne--Mumford 叠的拟紧拟分离态射 $f:\mathcal X\to\mathcal Y$, 定理~\bookref{5.5.16} 与推论~\bookref{5.5.17} 的相对版本也成立: $f$ 仿射, 当且仅当 $f_*:\QCoh(\mathcal X)\to\QCoh(\mathcal Y)$ 正合, 当且仅当对所有 $i>0$ 和 $F\in\QCoh(\mathcal X)$ 都有 $R^if_*F=0$.
\end{remark}
\begin{proposition}\label{src:5.5.19}
设 $X$ 是 Noether 代数空间. 若 $X_{\mathrm{red}}$ 是概形 (分别为拟仿射, 仿射), 则 $X$ 也是.
\end{proposition}
\begin{proof}
若 $X_{\mathrm{red}}$ 仿射, 则与~\cite[习题 III.3.1]{Har77} 完全一样, 用推论~\bookref{5.5.17} 说明 $X$ 仿射: 若 $F$ 是 $X$ 上的凝聚层, $I\subseteq\mathcal O_X$ 是定义 $X_{\mathrm{red}}$ 的幂零理想, 则利用滤过 $0=I^NF\subseteq I^{N-1}F\subseteq\cdots\subset IF\subseteq F$ 说明 $\mathrm H^i(X,F)$ 消失, 其逐次商 $I^kF/I^{k+1}F$ 的支集在 $X_{\mathrm{red}}$ 上. 若 $X_{\mathrm{red}}$ 是概形, 则每个点 $x\in|X|$ 都有开邻域 $U$, 使得 $U_{\mathrm{red}}$ 仿射. 因而 $U$ 仿射, $X$ 是概形. 若 $X_{\mathrm{red}}$ 拟仿射, 则 $X_{\mathrm{red}}\to\Spec\Gamma(X,\mathcal O_X)_{\mathrm{red}}$ 是开浸入. 因此 $X\to\Spec\Gamma(X,\mathcal O_X)$ 是开浸入, $X$ 拟仿射.
\end{proof}
\begin{theorem}[Chevalley 仿射性判据]\label{src:5.5.20}
设 $X\to Y$ 是 Noether 代数空间的有限满射态射. 若 $X$ 仿射, 则 $Y$ 也仿射.
\end{theorem}
\begin{proof}
可用推论~\bookref{5.5.17}, 仿照~\cite[习题 $4.1$]{Har77} 论证.
\end{proof}
\begin{exercise}\label{src:5.5.21}
证明下列推广: 若 $X\to Y$ 是代数空间的整满射态射, $Y$ 为 Noether, 且 $X$ 仿射, 则 $Y$ 仿射.

\textit{提示:} 将 $X=\lim_\lambda X_\lambda$ 写成有限 $Y$-概形的极限, 并使用极限方法 (\bookref{B.3}).
\end{exercise}
丰沛性的上同调判据~\cite[命题 $5.3$]{Har77} 也可推广到代数空间. 代数空间 $X$ 上的线丛 $L$ 称为丰沛的, 若 $X$ 是概形, 且 $L$ 在通常意义下丰沛.
\begin{exercise}\label{src:5.5.22}
设 $X$ 是 Noether 环上的固有代数空间. 对 $X$ 上的线丛 $L$, 证明 $L$ 丰沛, 当且仅当对 $X$ 上的每个凝聚层 $F$, 存在整数 $n_0$, 使得对 $i>0$ 和 $n\geq n_0$, 都有 $\mathrm H^i(X,F\otimes L^n)=0$. 另见~\cite[Tag 0D2W]{SP}.
\end{exercise}
与概形的情形一样 (命题~\bookref{B.2.9}), 丰沛性可在有限覆盖上检验.
\begin{exercise}\label{src:5.5.23}
设 $X$ 是 Noether 环上的固有代数空间, $L$ 是 $X$ 上的线丛. 若 $f:X'\to X$ 是有限满射态射, 则 $L$ 丰沛, 当且仅当 $f^*L$ 丰沛. 另见~\cite[Tag 0GFB]{SP}.
\end{exercise}
\begin{exercise}\label{src:5.5.24}
Noether 概形称为具有\newterm{Chevalley--Kleiman 性质}{Chevalley--Kleiman property}, 若每个有限点集都包含在某个仿射开集中. 证明若 $X\to Y$ 是 Noether 代数空间的有限满射态射, 且 $X$ 具有 Chevalley--Kleiman 性质, 则 $Y$ 也具有 Chevalley--Kleiman 性质. 见~\cite[推论 $48$]{Kol12}.
\end{exercise}

\subsection{沿域扩张的有效下降}\label{src:5.5.5-subsection}
\begin{lemma}\label{src:5.5.25}
设 $X$ 是域 $k$ 上局部有限型的拟分离代数空间. 若 $X_{\overline k}$ 是仿射概形, 则 $X$ 也是.
\end{lemma}
\begin{proof}
由于 $X_{\overline k}\to X$ 是整满射态射, 此陈述由 Chevalley 仿射性判据 (\bookref{5.5.21}) 推出.
\end{proof}
\begin{proposition}\label{src:5.5.26}
设 $X$ 是域 $k$ 上有限型的拟分离代数空间. 若 $X_{\overline k}$ 是概形, 则存在有限可分域扩张 $k\to K$, 使得 $X_K$ 是概形.
\end{proposition}
\begin{proof}
选取来自仿射概形的平展呈示 $U\to X$, 置 $R=U\times_X U$. 将 $\overline k=\operatorname{colim}k_\lambda$ 写成有限域扩张的余极限, 并置 $X_\lambda:=X_{k_\lambda}$, $U_\lambda=U_{k_\lambda}$, $R_\lambda=R_{k_\lambda}$. 设 $\overline V\subseteq X_{\overline k}$ 是仿射开子概形. 断言对 $\lambda\gg0$, 存在开子概形 $V_\lambda\subseteq X_\lambda$, 使得 $\overline V=U_\lambda\times_{k_\lambda}\overline k$. 事实上, $\overline V$ 在 $U_{\overline k}$ 中的原像 $\overline Z$ 的两个原像在 $R_{\overline k}$ 中相等. 由命题~\bookref{B.3.3} 和命题~\bookref{B.3.7}, 对 $\lambda\gg0$, 存在开子概形 $Z_\lambda\subseteq U_\lambda$, 使得 $\overline Z=Z_\lambda\times_{k_\lambda}\overline k$, 且 $Z_\lambda$ 在 $R_\lambda$ 中的两个原像相等. 由平展下降, $Z_\lambda$ 下降到所需的闭子概形 $V_\lambda\subseteq X_\lambda$. 引理~\bookref{5.5.25} 蕴含 $V_\lambda$ 是概形. 用有限多个仿射开集覆盖 $X_{\overline k}$, 并选取充分大的 $\lambda$, 得到 $k$ 的有限域扩张 $K=k_\lambda$, 使得 $X_\lambda$ 是概形. 现在将 $k\to k_\lambda$ 分解为可分域扩张 $k\to K$ 与纯不可分域扩张 $K\to k_\lambda$. 由于 $X_\lambda\to X_K$ 是有限万有同胚, Chevalley 仿射性定理 (\bookref{5.5.20}) 蕴含每个 $V_\lambda$ 在 $X_K$ 中的像仿射, 从而 $X_K$ 是概形.
\end{proof}
对 $X_{\overline k}$ 加上额外条件后, 可得 $X$ 是概形.
\begin{proposition}\label{src:5.5.27}
设 $X$ 是域 $k$ 上局部有限型的拟分离代数空间. 若 $X_{\overline k}$ 是概形, 且每个有限的 $\overline k$-点集都包含在某个仿射开集中 (例如 $X_{\overline k}$ 拟射影), 则 $X$ 是概形.
\end{proposition}
\begin{proof}
可假设 $X$ 拟紧. 将说明每个闭点 $x\in X$ 都有仿射开邻域. 设 $\Spec\kappa(x)\hookrightarrow X$ 是 $x$ 的剩余域的包含 (命题~\bookref{4.5.16}(3)), 并考虑 Cartesian 图
\[
\begin{tikzcd}
\coprod_{i=1}^n\Spec A_i\arrow[r,hook]\arrow[d]&X_{\overline k}\arrow[r]\arrow[d]&\Spec\overline k\arrow[d]\\
\Spec\kappa(x)\arrow[r,hook,"x"']&X\arrow[r]&\Spec k,
\end{tikzcd}
\]
其中 $n$ 是 $k$ 在 $K$ 中的可分闭包的次数, 每个 $A_i$ 是 Artin 局部 $\overline k$-代数. 对 $X_{\overline k}$ 的假设保证存在仿射开子概形 $\overline U\subseteq X_{\overline k}$, 包含每个 $\Spec A_i$ 的像. 由命题~\bookref{5.5.26}, 存在有限域扩张 $k\to K$, 使得 $X_K$ 是概形. 扩大 $K$ 后, 可使 $\overline U$ 下降为仿射开子概形 $U'\subseteq U_K$: 用命题~\bookref{B.3.3} 下降态射 $\overline U\to X$, 用命题~\bookref{B.3.7} 使其成为开浸入, 并用命题~\bookref{B.3.5} 保证仿射性. 注意 $U'$ 包含 $x$ 在 $X_K\to X$ 下的所有原像. 取 $K$ 的正规闭包, 可假设 $K$ 在 $k$ 上正规. 置 $G=\operatorname{Aut}(K/k)$, 则 $K^G$ 是 $k$ 的纯不可分域扩张. 于是 $G$ 在 $X_K$ 上自由作用, 且 $X_K/G=X_{K^G}$. $U'$ 在 $G$ 中各元素下的平移的交是 $G$-不变拟仿射簇 $U''$. 在 $U''$ 中选取包含 $x$ 的所有原像的仿射开子概形, 再对 $G$ 的平移取交, 得到包含 $x$ 的原像的 $G$-不变仿射开集 $V\subseteq X_K$. 则商 $V/G$ 是 $X_{K^G}$ 的仿射子概形, 包含 $x$ 的唯一原像 (定理~\bookref{5.2.3}). 设 $W$ 是 $V/G$ 在有限万有同胚 $X_{K^G}\to X$ 下的像, Chevalley 仿射性判据 (定理~\bookref{5.5.20}) 蕴含 $W$ 是 $x$ 的仿射邻域.
\end{proof}

\subsection{群代数空间是概形}\label{src:5.5.6-subsection}
域 $k$ 上的每个拟分离群代数空间都是概形. 当 $k$ 代数闭时, 这容易由定理~\bookref{5.5.1} 推出, 因为存在一个是概形的稠密开集, 可用有理点将其平移到各处. 一般情形依赖命题~\bookref{5.5.27}.
\begin{theorem}\label{src:5.5.28}
域 $k$ 上局部有限型的拟分离群代数空间 $G$ 是分离概形. 单位元的连通分支 $G^0$ 拟射影.
\end{theorem}
\begin{remark}\label{src:5.5.29}
若 $G$ 不拟分离, 上述推论不成立, 例如 $k$ 上的 $G=\mathbb G_a/\mathbb Z$ (例~\bookref{4.9.36}). 注意命题~\bookref{5.5.19} 蕴含结果在 Artin 基底上也成立. 在一般基概形上, 即使分离的群代数空间也未必是概形; 见~\cite[引理 X.14]{Ray70}.
\end{remark}
\begin{proof}[定理~\bookref{5.5.28} 的证明]
只需说明 $G$ 是概形. 事实上, 域上局部有限型的群概形必分离 (\bookref{B.1.19}(1)), 单位分支 $G^0$ 拟紧 (\bookref{B.1.19}(4)), 因而拟射影 (\bookref{B.1.19}(7)). 先假设 $k$ 代数闭. $G$ 有非空开子概形 $U$ (定理~\bookref{5.5.1}), 其中有一点 $h\in U(k)$. 对每个 $g\in G(k)$, 左乘 $gh^{-1}$ 定义同构 $G\xrightarrow{\sim}G$, $U$ 的像 $gh^{-1}U$ 是包含 $g$ 的概形. 一般情形由命题~\bookref{5.5.27} 推出, 其中利用了 $G_{\overline k}$ 是概形, 且每个有限点集都包含在某个仿射开集中 (引理~\bookref{5.5.30}). 另见~\cite[引理 $4.2$]{Art69b} 和~\cite[Tag 0B8D]{SP}.
\end{proof}
\begin{lemma}\label{src:5.5.30}
代数闭域 $k$ 上局部有限型的每个群概形 $G$ 都有以下性质: 每个有限的 $k$-点集都包含在某个仿射开子概形中.
\end{lemma}
\begin{proof}
设 $g_1,\ldots,g_n\in G(k)$. 先对 $n$ 归纳, 归结到所有元素 $g_i$ 都在同一连通分支的情形. 否则, 可写 $G=W_1\amalg W_2$, 其中 $W_1$ 有 $r$ 个点, $W_2$ 有 $n-r$ 个点, $0<r<n$. 由归纳, 存在仿射开集 $U_1\subseteq W_1$ 与 $U_2\subseteq W_2$, 分别包含这 $r$ 个点和 $n-r$ 个点. 则 $U_1\amalg U_2$ 是包含每个 $g_i$ 的仿射开集. 用 $g_1^{-1}$ 平移后, 还可假设 $g_1,\ldots,g_n\in G^0(k)$. 设 $U\subseteq G^0$ 是单位元的仿射开邻域. 由于 $G^0$ 不可约 (\bookref{B.1.19}(4)), $Ug_1^{-1}\cap\cdots\cap Ug_n^{-1}$ 非空, 且包含一个闭点 $h$. 由于 $h\in Ug_i^{-1}$, 每个 $g_i$ 都包含在仿射开集 $h^{-1}U$ 中.
\end{proof}
\begin{corollary}\label{src:5.5.31}
设 $\mathcal X$ 是对角线拟分离的代数叠. 则每个域值点的稳定子都是局部有限型的群概形.
\end{corollary}
\begin{proof}
由习题~\bookref{4.3.5}, $\mathcal X$ 的对角线局部有限型. 稳定子是对角线的基变换, 因而此陈述由定理~\bookref{5.5.28} 推出.
\end{proof}

\subsection{分离的一维代数空间是概形}\label{src:5.5.7-subsection}
\begin{theorem}\label{src:5.5.32}
满足 $\dim X\leq1$ 的分离 Noether 代数空间 $X$ 是概形.
\end{theorem}
\begin{proof}
可假设 $X$ 纯维数为 $1$. 由命题~\bookref{5.5.19}, 可假设 $X$ 约化. 由定理~\bookref{5.5.1}, 存在正则的稠密开子概形 $U\subseteq X$. 设 $D\subseteq U$ 是约化点的并, 使得 $X$ 的每个不可约分支都包含至少四个点. 则 $\mathcal O_U(D)$ 定义 $U$ 上的线丛, 可延拓为线丛 $\mathcal O_X(D)$. 由下一节证明的 Gabber 引理 (\bookref{5.6.1}), 存在来自概形 $Z$ 的有限满射 $f:Z\to X$. 将 $Z$ 换成其正规化后, 可假设 $Z$ 正则. 由于 $Z$ 是维数为 $1$ 的分离概形, $Z$ 拟射影, 且 $f^*\mathcal O_X(D)$ 丰沛. 由习题~\bookref{5.5.23}, $\mathcal O_X(D)$ 也丰沛. 另见~\cite[Tags 0ADD 和 09NZ]{SP}.
\end{proof}
\newpage
\begin{exercise}\label{src:5.5.33}
\begin{enumerate}[(a)]
\item (困难) 证明域上有限型, 维数为 $2$ 的每个光滑分离代数空间都是概形.
\item (中等) 构造一个域上有限型, 维数为 $2$ 的非正规分离代数空间, 它不是概形.
\item (开放) 每个维数为 $2$ 的分离正则代数空间都是概形吗?
\end{enumerate}
\end{exercise}
